\documentclass[conference]{IEEEtran}
\IEEEoverridecommandlockouts
% The preceding line is only needed to identify funding in the first footnote. If that is unneeded, please comment it out.
\usepackage{cite}
\usepackage{amsmath,amssymb,amsfonts}
\usepackage{algorithmic}
\usepackage{graphicx}
\usepackage{textcomp}
\usepackage{xcolor}
\usepackage{graphicx}
\usepackage{float}
\usepackage{framed}
\usepackage{enumerate}
\usepackage{minted}

\def\BibTeX{{\rm B\kern-.05em{\sc i\kern-.025em b}\kern-.08em
    T\kern-.1667em\lower.7ex\hbox{E}\kern-.125emX}}

\usepackage[english]{babel}
\usepackage[autostyle, english = american]{csquotes}
\MakeOuterQuote{"}

\begin{document}

\title{Contextware: Semantic Memory for Autonomous Software Development Agents}

\author{\IEEEauthorblockN{Justin Frauenhofer}
\IEEEauthorblockA{\textit{Computer Science and Engineering} \\
\textit{UC Santa Cruz}}
\and
\IEEEauthorblockN{Graeme Maclean}
\IEEEauthorblockA{\textit{Computer Science and Engineering} \\
\textit{UC Santa Cruz}}
}

\maketitle

\begin{abstract}
The integration of Large Language Models (LLMs) into software development workflows has significantly accelerated code generation and debugging; however, these systems frequently suffer from "contextual amnesia", the inability to retain architectural decisions, user preferences, and high-level design intent across disconnected development sessions. This paper presents Contextware, a persistent semantic memory skill developed for the Gemini CLI that addresses these limitations through a localized, three-tier memory architecture.

Contextware utilizes a local vector database (LanceDB) and CPU-efficient embeddings (FastEmbed) to manage three distinct data collections: a "Hippocampus" for explicit developer facts and constraints, a "Diary" for historical conversation logs, and a "Map" for a semantically indexed codebase. A novel contribution of this work is the implementation of an autonomous background indexing mechanism that employs "headless" agent instances to recursively summarize file purposes and symbol structures without interrupting the primary developer workflow. By combining proactive memory capture with a robust Retrieval-Augmented Generation (RAG) pipeline, Contextware enables CLI agents to maintain a consistent project "mental model," significantly reducing the need for manual context re-provisioning. Preliminary results demonstrate that this architecture improves agent autonomy in complex refactoring tasks and ensures long-term adherence to project-specific conventions, providing a scalable blueprint for more reliable autonomous engineering systems.


Keywords: Semantic Memory, Autonomous Agents, Software Engineering, Retrieval-Augmented Generation (RAG), Large Language Models, Vector Databases.
\end{abstract}

\section{Introduction}

The emergence of Large Language Models (LLMs) has catalyzed a paradigmatic shift in software engineering, transitioning from syntax-level autocompletion to the deployment of \textbf{autonomous agents} capable of resolving real-world GitHub issues. As the field moves from simple code generation to \textbf{intent abstraction} [MemoryManagementLCNC], agents are increasingly expected to act as "deep agents" that automate complex, long-duration business processes. However, as these agents move from statement-level generation to repository-level problem solving, they encounter a fundamental bottleneck: \textbf{agent amnesia}.

\subsection{The Context Window Bottleneck}
Autonomous coding agents are currently restricted by the \textbf{fixed context windows} of Transformer architectures. While recent advancements have expanded these limits to millions of tokens, the self-attention mechanism's \textbf{quadratic complexity $O(n^2)$} [Nemori] makes simply increasing the context window prohibitively expensive in terms of latency and computational cost. In professional software development, interactions with a codebase can span weeks, months, or even indefinitely. Holding the entire historical context within a single prompt is not only computationally unsustainable but also introduces significant noise that hinders the agent's ability to focus on task-relevant information.

\subsection{The Engineering Gap and Multi-Hop Reasoning}
Professional coding is uniquely memory-intensive. Real-world repositories, such as those represented in the \textbf{SWE-bench} framework, average 438,000 lines of code across thousands of files [SWE-bench]. Resolving a single bug often requires \textbf{multi-hop reasoning}---the ability to synthesize information dispersed across multiple functions, classes, and modules defined in disparate development sessions. Standard models often suffer from the \textbf{``lost in the middle'' phenomenon} [SWE-bench, Nemori], where attention degrades for information buried in the center of a long prompt. This manifests as \textbf{memory inflation}: the linear accumulation of raw logs increases token costs and latency while simultaneously introducing irrelevant context that leads to self-degradation and error propagation [MemoryManagementLCNC].

\subsection{Operational and Economic Bottlenecks}
The failure to manage memory is not merely a technical limitation but an operational one. Without a persistent memory mechanism, agents are forced to ``re-learn'' repository structures and developer preferences in every new session. This leads to a state of \textbf{contextual degradation}, where the accumulation of flawed or irrelevant memories causes the agent's performance to decline over time [MemoryManagementLCNC]. Furthermore, the economic cost of re-sending hundreds of thousands of tokens of repository context for every sub-task request makes autonomous agents impractical for large-scale enterprise deployment.

\subsection{Contributions of Contextware}
To bridge this gap, we introduce \textbf{Contextware}, a memory-augmented skill for the Gemini CLI that shifts the agent from passive logging to \textbf{proactive, tiered memory management}. The specific contributions of this work include:
\begin{itemize}
    \item A \textbf{three-tier storage hierarchy} (Hippocampus, Diary, and Map) inspired by the Common Model of Cognition to segregate persistent constraints from historical task logs.
    \item An \textbf{autonomous background indexing engine} that utilizes headless agent instances to recursively summarize file intent and repository structure without interrupting the primary user workflow.
    \item A \textbf{structural symbol navigation} module that uses Abstract Syntax Tree (AST) parsing to extract class signatures and method maps, ensuring the semantic integrity of retrieved code.
    \item A design optimized for \textbf{local deployment} using LanceDB and FastEmbed, ensuring developer privacy and millisecond-level retrieval latency.
\end{itemize}

\section{Background}

The development of Contextware is rooted in a synthesis of cognitive science, systems engineering, and natural language processing to solve the foundational challenge of agent amnesia established in the previous section. This section provides an exposition of the cognitive frameworks used to model persistent intelligence, the operational lifecycle of agentic memory, and the technical limitations of existing retrieval systems when applied to complex software engineering tasks.

\subsection{Cognitive Memory Frameworks}
To simulate human-like persistence and overcome the amnesia established in our Introduction, researchers have turned to the \textbf{Common Model of Cognition (CMC)} [MemoryMatters]. The CMC posits that a functional intelligent agent must segregate information based on its temporal and conceptual nature rather than treating interaction logs as a monolithic stream [SurveyAgentMemory]. Contextware adopts this tripartite approach to memory storage:
\begin{itemize}
    \item \textbf{Working Memory (WM)}: Functionally represented by the immediate context window, this tier is populated with the situation description required to drive the agent's current prompt [MemoryMatters].
    \item \textbf{Semantic Memory (SM)}: Represents accumulated factual knowledge and abstract concepts that remain independent of specific events [A-Mem]. In the coding domain, this includes project-wide architectural constraints and developer-specific preferences.
    \item \textbf{Episodic Memory (EM)}: A narrative log of past temporal experiences, recording ``what happened when'' [MIRIX]. This allows the agent to recall specific past successes and, crucially, avoid re-attempting failed sub-tasks [MemoryManagementLCNC].
    \item \textbf{Procedural Memory (PM)}: Stores actionable skills and structural maps of the environment [CodeAgent]. For Contextware, this is realized as a semantically indexed map of the codebase.
\end{itemize}

\subsection{From Static Retrieval to Dynamic Experience}
The transition from Retrieval-Augmented Generation (RAG) to \textbf{Memory-Augmented Generation (MAG)} represents a fundamental shift in how LLMs interact with data [Nemori]. Traditional RAG systems were designed for static knowledge bases and operate on a \textbf{``stateless information patch''} philosophy [Nemori]. While effective for finding localized facts, these systems are misaligned with the needs of an autonomous agent for three reasons: (1) their reliance on offline indexing prevents real-time learning; (2) their focus on document similarity ignores the temporal nature of dialogue; and (3) their use of \textbf{arbitrary fixed-size chunking} often leads to \textbf{semantic fragmentation} [Nemori]. 

In a software repository, fragmentation occurs when a single class or method is split across multiple vector chunks, destroying the agent's ability to perceive the logical integrity of the code. MAG architectures, by contrast, emphasize the \textbf{autonomous self-organization} of an agent's own lived experience into optimized, coherent representations [Nemori, MemInsight].

\subsection{The Operational Lifecycle: Writing, Management, and Reading}
An effective persistent memory must facilitate a complete lifecycle of operations to prevent the knowledge base from becoming a haphazard pile of raw text [SurveyAgentMemory, Mem0].
\begin{enumerate}
    \item \textbf{Memory Writing}: This is the process of \textbf{proactive information extraction}. Rather than passively logging every token, the agent must identify and distill salient facts and event summaries for storage [Mem0, Contextware].
    \item \textbf{Memory Management}: This involves resolving conflicts and maintaining consistency. As will be discussed in the Relevant Works section regarding Mem0, this requires operations such as \texttt{ADD} for new information, \texttt{UPDATE} for augmenting existing knowledge, and \texttt{DELETE} for removing contradicted facts [Mem0, Algorithm 1].
    \item \textbf{Memory Reading}: This utilizes \textbf{active retrieval} to prevent the model from defaulting to its internal parametric knowledge, which may be outdated or incorrect [MIRIX, ActiveRetrievalVisual].
\end{enumerate}

\subsection{Scalability and the Logic of Forgetting}
As an interaction history grows over weeks or months, agents suffer from \textbf{memory inflation}, where the accumulation of irrelevant or redundant data leads to exponential increases in token costs and latency [MemoryManagementLCNC]. Without a management mechanism, this leads to \textbf{contextual degradation}, a state where the agent begins to contradict its own previous decisions or repeat errors due to ``misaligned experience replay'' [MemoryManagementLCNC].

To ensure lifelong learning, agents must implement \textbf{intelligent forgetting} or \textbf{intelligent decay} [MemoryOS, MemoryManagementLCNC]. Inspired by the biological principle that memory utility is a function of engagement, systems assign a \textbf{Utility Score} $S(M_i)$ to each memory entry. This score is typically a composite of \textbf{Recency} (time-decay), \textbf{Relevance} (semantic similarity), and \textbf{Frequency} (access counts) [MemoryOS, Algorithm 1]. High-utility ``hot'' memories are retained in faster tiers, while ``cold,'' low-utility memories are consolidated into compact summaries or evicted entirely to maintain system efficiency.

\subsection{Memory Requirements for Repository-Level Reasoning}
The complexity of repository-level software engineering tasks, such as those represented in the \textbf{SWE-bench} framework, places unique demands on an agent's memory [SWE-bench]. Success in this domain hinges on \textbf{multi-hop reasoning}---the ability to synthesize information dispersed across thousands of files and disparate development sessions [SWE-bench, AriGraph]. 

For example, an agent must bridge the gap between a high-level requirement defined in session A and a specific utility function implemented in session B to solve a bug in session C. Standard vector retrieval often suffers from \textbf{relational blindness} in these scenarios. This gap necessitates the \textbf{structural symbol navigation} and \textbf{recursive background indexing} approaches that Contextware implements to maintain a global ``mental model'' of the project without overloading the immediate context window.

\section{Relevant Works}
Current research has transitioned from linear text logs to structured and proactive memory management architectures.

\subsection{Structural Paradigms: Graphs and Hierarchies}
Several recent works emphasize the importance of structured representation for complex reasoning. AriGraph integrates semantic factual triplets and episodic experience edges within a world model graph, significantly improving spatial orientation and multi-hop reasoning in text-based environments [AriGraph]. Mem0 and its variant Mem0g utilize graph-based representations to capture complex relational structures, achieving superior accuracy in temporal and open-domain tasks compared to unstructured memory [Mem0]. A-MEM draws inspiration from the Zettelkasten method, using interconnected atomic notes that evolve over time as new memories are integrated, allowing the system to refine its contextual understanding [A-Mem]. While these \textbf{structured memory techniques} provide significant performance gains for particular workflows, their complexity can lead to higher latency and error accumulation during graph construction [Mem0].

\subsection{Operating System-Inspired Architectures}
Other architectures utilize principles from computer operating systems. \textbf{MemoryOS} introduces a three-tier storage hierarchy comprising short-term, mid-term, and long-term personal memory. It employs \textbf{segmented paging} and \textbf{heat-based eviction} strategies based on access frequency and recency, demonstrating that mid-term topic summaries have the highest impact on maintaining dialogue coherence [MemoryOS]. \textbf{MemGPT} also treats memory like an operating system with a ''RAM/Disk'' analogy, using self-directed function calls to page information in an d out of the limited context window. These hierarchical tiers ensure that topic-aligned content remains accessible while maintaining consistency with user-specific preferences [MemGPT].

\subsection{Proactive Learning and Multimodal Integration}
Moving beyond passive extraction, Nemori introduces the Predict-Calibrate Principle, where the agent forecasts its own experience and only distills new knowledge from the surprise or prediction gap against the raw interaction [Nemori]. MemInsight autonomously mines and prioritizes semantic attributes like user intent and emotion, which significantly improves recall compared to raw text similarity [MemInsight]. For multimodal contexts, MIRIX manages memory across six specialized components, including a secure Knowledge Vault for verbatim sensitive data and a Resource Memory for full documents, and can process thousands of screenshots to build a visual activity history [MIRIX].

\subsection{Specialized Coding Agents and Multi-Agent Teams} In the domain of software engineering, specialized agents leverage memory to handle repo-level tasks. \textbf{CodeAgent} formalizes repository-level code generation by integrating tools for \textbf{symbol navigation} and documentation reading into its agentic workflow [CodeAgent]. \textbf{RepairAgent} is the first autonomous agent specifically for \textbf{program repair}, using a \textbf{finite state machine} to guide the model through gathering information and searching for repair ingredients [RepairAgent].
Furthermore, the research community is increasingly interested in \textbf{multi-agent coding teams}. \textbf{MetaGPT} incorporates \textbf{Standardized Operating Procedures (SOPs)} into agent collaboration by encoding role specialization (e.g., Architect, Engineer, QA) into structured prompt sequences. \textbf{ChatDev} utilizes a \textbf{``chat chain''} workflow where specialized agents engage in multi-turn dialogues to autonomously craft solutions, effectively reducing coding hallucinations [MetaGPT and ChatDev]. While these multi-agent environments are a high-interest area for future investigation, the current implementation of \textbf{Contextware} focuses on a single-agent architecture to minimize synchronization overhead and manage complexity within the project's scope.

\subsection{Key Design Takeaways for the Project}
The following insights from the research are critical for implementing an efficient agent memory:

\begin{itemize}
    \item \textbf{Segregation of Memory Types}: Procedural, episodic, and semantic memories should be stored in separate structures to prevent internal contradictions and improve retrieval precision.
    Intelligent Decay Mechanism: A utility score based on recency, relevance, and user feedback must be implemented to prune irrelevant information and avoid performance decline caused by catastrophic interference.
    \item \textbf{Active Retrieval}: Instead of waiting for explicit triggers, the agent should self-generate a current topic to search memory proactively, preventing it from defaulting to outdated internal knowledge.
    \item \textbf{Human-in-the-Loop Transparency}: Providing a visual interface where users can pin important facts or delete flawed memories enhances reliability and user trust, particularly in business processes.
    \item \textbf{Semantic Attribute Mining}: Annotating raw data with autonomously mined attributes is more effective than relying solely on raw text embeddings for retrieval.
\end{itemize}

\section{Integration of Memory Architecture Insights}

Our project, \textbf{Contextware}, is explicitly designed to address key limitations in current LLM agent memory management as identified by Hatalis et al. \cite{Hatalis2024Memory}. The article highlights several critical challenges in the development of autonomous agents that we have directly incorporated into our architectural design.

\subsection{Separation of Memory Types}
Hatalis et al. identify the lack of ``separation between types of memories'' as a fundamental problem, noting that agents often fail when procedural, episodic, and semantic memories are conflated in a single vector space \cite{Hatalis2024Memory}. To resolve this, Contextware implements three distinct collections within our local vector database:

\begin{itemize}
    \item \texttt{facts} (Semantic Memory): Stores explicit user preferences and constraints.
    \item \texttt{episodes} (Episodic Memory): Records the outcomes of specific tasks.
    \item \texttt{code\_index} (Procedural Memory): Maps the codebase structure and symbols.
\end{itemize}

This distinct separation ensures that a user preference stored in \texttt{facts} does not compete with a past debugging attempt stored in \texttt{episodes} during vector retrieval, preventing the ``poor retrieval'' issues cited in the article.

\subsection{Mitigating Infinite Subtask Loops}
A significant reliability issue raised in the study is the ``Infinite subtask loop,'' where agents repeatedly fail at the same task because they do not ``detect if failures are reoccurring'' \cite{Hatalis2024Memory}. 

We address this by introducing the \texttt{episodes} collection, which specifically records the \texttt{goal}, \texttt{result} (success/failure), and a \texttt{summary} of actions taken. The retrieval workflow in \texttt{recall.py} is designed to query this ``Black Box'' before the agent attempts a logical task. This allows the agent to recognize past failures and adjust its strategy rather than repeating the same mistake, directly answering the article's call for memory systems that facilitate adaptation.

\subsection{Scalable Procedural Memory}
The article discusses the necessity for agents to leverage external knowledge sources and metadata to enhance code generation \cite{Hatalis2024Memory}. 

Contextware solves the problem of fixed-size context windows by using background ``headless'' Gemini instances to generate summaries and symbol maps for the \texttt{code\_index}. This creates a semantic map of the project, allowing the agent to perform broad searches over the entire file structure without loading every file into its immediate context.

\subsection{Lifelong Memory Management}
Finally, \cite{Hatalis2024Memory} notes that agents accumulate extensive memories over time, leading to issues with ``indexing, scaling, and metadata-filtered queries.'' 

To ensure sustainable performance, we utilize \texttt{lancedb} for high-performance local vector storage and implement strict metadata tracking. Specifically, the \texttt{code\_index} tracks the \texttt{mtime} (last modified timestamp) of indexed files. This metadata-driven approach ensures the agent can distinguish between ``stale'' and ``fresh'' information, addressing the article's requirement for memory systems that maintain reactivity to environmental changes.

\section{Implementation}
Contextware is implemented as a Python-based skill for the Gemini CLI, using \texttt{uv} to execute scripts with PEP 723 inline dependency metadata to ensure isolated and reproducible execution without a traditional virtual environment. The architecture leverages local, high-performance components to minimize latency and ensure data privacy.

\subsection{Vector Storage and Embeddings}
We utilize \texttt{LanceDB}, a serverless, persistent vector database designed for high-performance random access. Unlike traditional vector stores that require a running service, \texttt{LanceDB} stores data in the Lance columnar format, allowing for efficient disk-based retrieval. For generating semantic representation, we employ the \texttt{FastEmbed} library with teh \texttt{BAAI/bge-small-en-v1.5} model. This model produces 384-dimensional vectors and is optimized for CPU execution, making it ideal for background indexing tasks on developer machines without GPU acceleration.

\subsection{Data Modeling and Schema}
Data structures are defined using \texttt{Pydantic} and \texttt{LanceModel} to ensure strict typing and seamless integration with the database. We maintain three primary collections:

\begin{itemize}
    \item \textbf{Facts}: Stores developer preferences using the \texttt{Fact} model, which includes the raw \texttt{content}, a \texttt{timestamp}, and the corresponding \texttt{vector}.
    \item \textbf{Episodes}: Captures task outcomes via the \texttt{Episode} model. It includes the \texttt{goal}, \texttt{result} (success, failure, or partial), \texttt{category}, and a narrative \texttt{summary}. The embedding vector is generated from a concatenation of the goal and summary to capture both intent and action.
    \item \textbf{Code Index}: Maps the codebase structure using the \texttt{CodeIndex} model. It stores the absolute \texttt{file\_path}, a semantic \texttt{summary}, and hierarchical metadata including a JSON-serialized map of \texttt{classes} to their methods and a list of \texttt{top\_level\_functions}. The \texttt{last\_modified} timestamp is recorded to enable staleness detection.
\end{itemize}

\subsection{CLI Interface and Workflows}
The agent interacts with Contextware through two primary scripts: \texttt{store.py} for data ingestion and \texttt{recall.py} for retrieval. 

The \texttt{store.py} script includes an automated symbol extraction engine for Python files. Using the \texttt{ast} (Abstract Syntax Trees) module, it recursively parses source code to identify class definitions and function signatures. This allows the agent to index complex files with minimal manual input.

The \texttt{recall.py} script provides a multi-scope search interface. It supports semantic similarity search across all collections or specific sub-tiers (e.g., searching only \texttt{facts}). Additionally, it implements a direct path-lookup mechanism that compares the stored \texttt{last\_modified} timestamp with the current file system \texttt{mtime}. If a discrepancy is detected, the result is flagged as \texttt{[STALE]}, signaling to the agent that a re-indexing operation is required.

\clearpage
\bibliographystyle{IEEEtran}
\bibliography{references}


\end{document}
